%20 min preso!
\documentclass[xcolor=table]{beamer}
\usepackage{beamerthemesplit}
\usepackage{wrapfig}
\usetheme{SPbGU}
\usepackage{pdfpages}
\usepackage{amsmath}
\usepackage{cmap} 
\usepackage[T2A]{fontenc} 
\usepackage[utf8]{inputenc}
\usepackage[english]{babel}
\usepackage{indentfirst}
\usepackage{amsmath}
\usepackage{tikz}
\usepackage{multirow}
\usepackage[noend]{algpseudocode}
\usepackage{algorithm}
\usepackage{algorithmicx}
\usepackage{fancyvrb}
\usetikzlibrary{shapes,arrows}

\usepackage{tabularx}
%\newcolumntype{P}[1]{>{\centering\arraybackslash}p{#1}}
%\newcolumntype{M}[1]{>{\centering\arraybackslash}m{#1}}
\newcolumntype{Y}{>{\raggedleft\arraybackslash}X}


%usepackage{fancyvrb}
%\usepackage{minted}
%\usepackage{verbments}

\newtheorem{mytheorem}{Theorem}
\renewcommand{\thealgorithm}{}


\beamertemplatenavigationsymbolsempty

\title[DNN + Formal Grammars]{The Composition of Dense Neural Networks and Formal Grammars for Secondary Structure Analysis}
%\subtitle[YaccConstructor]{Parsing techniques for graph analysis}
% То, что в квадратных скобках, отображается в левом нижнем углу. 
\institute[JetBrains Research]{
JetBrains Research, Programming Languages and Tools Lab  \\
Saint Petersburg University
}

% То, что в квадратных скобках, отображается в левом нижнем углу.
\author[Semyon Grigorev]{Polina Lunina, \textbf{Semyon Grigorev}}

\date{Febrary !!!, 2019}

\definecolor{orange}{RGB}{179,36,31}

\begin{document}
{
\begin{frame}[fragile]
  \begin{table}
  \centering
  \begin{tabularx}{\linewidth}{YcX}
    \includegraphics[height=1.5cm]{pictures/jetbrainsResearch.pdf} \hfill
    & \begin{minipage}[t]{0.3\textwidth}\center \vspace{-1cm}  BIOINFORMATICS 2019 
      \end{minipage}
    & \hfill \includegraphics[height=1.5cm]{pictures/SPbGU_Logo.png}
  \end{tabularx}
  \end{table}
  %\begin{tabular}{p{1.5cm} p{7.5cm} p{1cm}}
  % \begin{center}
  %    \includegraphics[height=1.5cm]{pictures/jetbrainsResearch.pdf}
  %  \end{center}
  %  &
  %  \begin{center}
  %    %\includegraphics[height=1.5cm]{pictures/icfplogo.png}
  %    BIOINFORMATICS \\ 2019
  %  \end{center}
  %  &
  %  \begin{center}
  %    \includegraphics[height=1.5cm]{pictures/SPbGU_Logo.png}
  %  \end{center} 
  %\end{tabular}
  \titlepage
\end{frame}
}

\begin{frame}[fragile]
  \transwipe[direction=90]
  \frametitle{Grammar}
\begin{verbatim}
s1: stem<s0>
any_str : any_smb*[2..10]
s0: any_str | any_str stem<s0> s0
any_smb: A | T | C | G
stem1<s>:               \\ stem of height exactly 1
      A s T | G s C | T s A | C s G 
stem2<s>:               \\ stem of height exactly 2 
      stem1< stem1<s> >
stem<s>:                \\ stem of height 3 or more                                  
      A stem<s> T 
    | T stem<s> A 
    | C stem<s> G 
    | G stem<s> C 
    | stem1< stem2<s> >  
\end{verbatim}

\end{frame}

\begin{frame}[fragile]
  \transwipe[direction=90]
  \frametitle{Example 1: Stem}
\centering
 $\omega_1=$ \texttt{CCCCATTGCCAAGGACCCCACCTTGGCAATCCC}
\vspace{1cm}

\includegraphics[width=.8\textwidth]{pictures/4.pdf}

\end{frame}

\begin{frame}[fragile]
  \transwipe[direction=90]
  \frametitle{Example 2: Pseudoknot}
\centering
 $\omega_2=$ \texttt{CCACTTACCTATGACCTAAGTCCTCATACC}
\vspace{1cm}

\includegraphics[width=.8\textwidth]{pictures/5.pdf}

\end{frame}


\begin{frame}[fragile]
  \transwipe[direction=90]
  \frametitle{Example 2: Pseudoknot}
\centering
 $\omega_3=$ \texttt{\\CAGGGCATAACCTAGCCCAACCTTGCCAAGG\\TTGGGGTCGAGGGTTCGAATCCCTTCGCCCGCTCCA}
\vspace{1cm}
\begin{minipage}[t]{0.48\textwidth}
\centering
\includegraphics[width=.94\textwidth]{pictures/Fold1.pdf}
\end{minipage}
\begin{minipage}[t]{0.48\textwidth}
\centering
\includegraphics[width=.94\textwidth]{pictures/Fold2.pdf}
\end{minipage}

\end{frame}

\begin{frame}[fragile]
  \transwipe[direction=90]
  \frametitle{Example 2: Pseudoknot}
\centering
%\begin{minipage}[t]{0.48\textwidth}
%\centering
%\includegraphics[width=.45\textwidth]{pictures/Fold1.pdf}
%\end{minipage}
%\begin{minipage}[t]{0.48\textwidth}
%\centering
%\includegraphics[width=.45\textwidth]{pictures/Fold2.pdf}
%\end{minipage}
\includegraphics[width=\textwidth]{pictures/0m.pdf}
\end{frame}

\begin{frame}
  \transwipe[direction=90]
  \frametitle{F\# type providers}
\begin{itemize}
 \item Compile-time metaprogramming for types creation
 \begin{itemize}
  \item Type provider is a \textbf{function which constructs type}
 \end{itemize}
 \item Design-time features in IDE
 \begin{itemize}
   \item Completion
   \item Type information
 \end{itemize}
 \item Used for type-safe integration of external data with fixed schema
 \begin{itemize}
   \item Type providers for XML, JSON, INI
   \item R, SQL
 \end{itemize}
\end{itemize}

\end{frame}


\begin{frame}[fragile]
  \transwipe[direction=90]
  \frametitle{Example of INI type provider}

\small{
\begin{verbatim}
[Section1]
intSetting = 2
stringSetting = stringValue
[Section2]
floatSetting = 1.23
boolSetting = true
anotherBoolSetting = False
emptySetting =
stringWithSemiColonValue = DataSource=foo@bar;UserName=blah 
\end{verbatim}
}
%\includegraphics[width=0.6\textwidth]{pictures/INIExample.png}

\end{frame}

\begin{frame}
  \transwipe[direction=90]
  \frametitle{OpenCL C type provider}
\begin{itemize}
\item We want to construct type-safe wrapper for existing library
\item OpenCL standard declares source-level distribution with in place compilation
\begin{itemize}
\item[+] We can work with source code, not with binaries
\item[--] Existing library is a set of files includes *.h files
\end{itemize}
\item It is enough to process functions signatures
\end{itemize}
\end{frame}
 
\begin{frame}
  \transwipe[direction=90]
  \frametitle{OpenCL C type provider: architecture}
    
 % \includegraphics[width=0.95\textwidth]{pictures/OpenCL_C_TP.pdf}
 % \\
  \pause
  \center{Yes, it is typical type provider}
\end{frame}

\begin{frame}
  \transwipe[direction=90]
  \frametitle{Limitations}
\begin{itemize}
\item Only (small) subset of OpenCL C
 \begin{itemize}
 \item *.h files are not supported
 \item preprocessor is not supported
 \item only small subset of syntax is supported
  \end{itemize}
\item Very simple C to F\# type mapping
\end{itemize}

\end{frame}

            
\begin{frame}
  \transwipe[direction=90]
  \frametitle{Examples}         
  %\includegraphics[width=0.95\textwidth]{pictures/smpl1.png}
  %\pause
  %\\
  %\includegraphics[width=0.95\textwidth]{pictures/smpl3.png}
  %\pause
  %\\
  %\includegraphics[width=0.95\textwidth]{pictures/smpl2.png}

\end{frame}     

\begin{frame}
  \transwipe[direction=90]
  \frametitle{Future work}         
\begin{itemize}
\item Improve OpenCL C support
\begin{itemize}
\item Lexer and parser
\item Translator 
\item Types mapping
\item Headers files processing
\item \dots
\end{itemize}
\item Unify kernels on client side
\begin{itemize}
\item Currently native Brahma.FSharp's kernel and kernel loaded by type provider are different 
types
\end{itemize}
\item Improve mechanism of kernels composition
\end{itemize}
\end{frame}     
            
            
\begin{frame}
  \transwipe[direction=90]
  \frametitle{Summary}         
\begin{itemize}
\item F\# OpenCL C type provider
\begin {itemize}
\item Type-safe integration of existing OpenCL C code in F\# applications
\item Proof of concept
\end{itemize}
\vspace{1cm}
\item Source code on GitHub: \url{https://github.com/YaccConstructor/Brahma.FSharp}
\item Package on NuGet: \url{https://www.nuget.org/packages/Brahma.FSharp/}
\end{itemize}
\end{frame}           
            
\begin{frame}
\transwipe[direction=90]
\frametitle{Contact Information}
\begin{itemize}
  \item Semyon Grigorev: \href{mailto:s.v.grigoriev@spbu.ru}{s.v.grigoriev@spbu.ru}
  \item Kirill Smirenko: \href{mailto:k.smirenko@gmail.com}{k.smirenko@gmail.com}
\end{itemize}
\begin{itemize}
  \item Brahma.FSharp: \href{https://github.com/YaccConstructor/Brahma.FSharp}{https://github.com/YaccConstructor/Brahma.FSharp}
\end{itemize}
\hspace{2cm}
\center{\huge{Thanks!}}
\end{frame}
\end{document}
